\documentclass{article}

\begin{document}

\title{Generating Parallel Numerical Codes for PDE's
 \thanks{Research sponsored in part by the Phillips Laboratory, Air Force Materiel  Command, USAF, under cooperative agreement
number F29601-93-2-0001. Views and conclusions contained in this document are those of the author(s) and should not be
interpreted as necessarily representing the official policies or endorsements, either expressed or implied, of Phillips Laboratory
or the U.S. Government.; 
CCR-9531828 NSF; 
ME 050 (1997) program KONTAKT, Czech Republic.}
}

\author{Dobromil Pryl}

\date{}

\maketitle

\begin{center}
Faculty of Nuclear Sciences and Physical Engineering,
 Czech Technical University, Prague
\end{center}
E-mail: pryl@atlas.cz

\begin{abstract}
The talk presents a package for generating numerical programs 
for solving PDE's using finite differences method.
The package is implemented in LISP programing language
in REDUCE computer algebra system, interfacing to many existing software tools.
Complete FORTRAN source codes can be generated,
including all declarations, output of results, and comments.
It is easy to switch between creating sequential and parallel codes.
Parallelization is done by domain decomposition using message passing, 
utilising the MPI standard.
Parallelism is supported for both implicit and explicit difference schemes.
Applications to several simple (one-way wave equation)
and more complex (shallow water equations) problems
are included for illustration and parallel speedup analysis.
\end{abstract}

\end{document}
